\subsection{理想}

由恒等式 (3) 可知，在李代数 \(\mathfrak{g}\) 中左理想与右理想没有区别，每个理想都是双边的。因此我们简称为理想。

*例. 设 G 是李群，g 是其李代数，H 是 G 的李子群。H 上每个左不变向量场都自然地定义 G 上的一个左不变向量场，从而 H 的李代数 h 自然嵌入 g；在此嵌入下，h 被视为 g 的李子代数。若 H 是 G 的正规子群，则 h 在 g 中的典范像是 g 的理想。*

g 的理想是 g 的一个在 g 的内导子下稳定的子模。

定义 4. g 的一个在 g 的每个导子下都稳定的子模称为 g 的特征理想。

命题 2. 设 g 是李代数，a 是 g 的理想（分别为特征理想），b 是 a 的特征理想。则 b 是 g 的理想（分别为特征理想）。

g 的每个内导子（分别为每个导子）都保持 a 稳定，并在 a 上诱导一个导子，因而保持 b 稳定。

设 \(\mathfrak{g}\) 是李代数。若 \(\mathfrak{a}\) 和 \(\mathfrak{b}\) 是 \(\mathfrak{g}\) 的理想，则 \(\mathfrak{a} + \mathfrak{b}\) 和 \(\mathfrak{a} \cap \mathfrak{b}\) 也是 \(\mathfrak{g}\) 的理想。

设 \(\mathfrak{a}\) 和 \(\mathfrak{b}\) 是 \(\mathfrak{g}\) 的两个子模。滥用记号，将 \(\mathfrak{g}\) 中由形如 \([x,y](x\in\mathfrak{a},y\in\mathfrak{b})\) 的元素生成的子模记作 \([\mathfrak{a},\mathfrak{b}]\)。由恒等式 (3)，有 \([\mathfrak{a},\mathfrak{b}] = [\mathfrak{b},\mathfrak{a}]\)。若 \(z\in\mathfrak{g}\)，则 \([z,\mathfrak{a}]\) 或 \([\mathfrak{a},z]\) 表示子模 \([\mathrm{K}z,\mathfrak{a}] = (\mathrm{ad}\,z)(\mathfrak{a})\)。

命题 3. 若 a 和 b 是 g 的理想（分别为特征理想），则 {[}a, b{]} 是 g 的理想（分别为特征理想）。

设 D 是 g 的内导子（分别为导子）。若 \(x \in a\) 且 \(y \in b\)，则

\[
\mathrm{D} ([ x, y ]) = [ \mathrm{D} x, y ] + [ x, \mathrm{D} y ] \in [ \mathfrak {a}, \mathfrak {b} ].
\]

故命题得证。

若 \(\mathfrak{a}\) 是 \(\mathfrak{g}\) 的子模，则由 \(x \in \mathfrak{g}\) 且满足 (ad \(x\) ). \(\mathfrak{a} \subset \mathfrak{a}\) 的元素组成的集合是一个子代数 \(\mathfrak{n}\)，属于 \(\mathfrak{g}\)，称为 \(\mathfrak{a}\) 在 \(\mathfrak{g}\) 中的正规化子。若进一步 \(\mathfrak{a}\) 是 \(\mathfrak{g}\) 的子代数，则 \(\mathfrak{a} \subset \mathfrak{n}\)，且 \(\mathfrak{a}\) 是 \(\mathfrak{n}\) 的理想。

